\subsection{All local sections of the evaluated elliptic quotient}

\begin{statement}{E.60}{Proved}{Complete local coordinates and the exact integral class-three image at two}
Let $R$, $M$, $\xi,\eta,U,V$ and the genuine Frobenius
$\varphi$ be those of E.58. Work over the actual local field
$K_v=\mathbb Q_2(\zeta_3)$, with residue field $\mathbb F_4$,
at the prime of $K=\mathbb Q(\zeta_3)$ above two. Use the
same actual prolongation and discriminant-root convention.
All geometric-conjugacy classes of local sections of $R$
are the following explicit disjoint union:
\[
 \boxed{\quad
 \coprod_{A,B\in\{0,1,2\}}
       \mathcal L_{A,B}\times
              \bigl((\mathbb Z/3)^7\oplus(\mathbb Z/27)^2\bigr),
 \quad}
 \tag{E60.1}
\]
where
\[
 \begin{gathered}
 h=(h_{ij})_{0\le i,j\le2}\in\mathbb Z_3^9,
 \qquad h_0=\sum h_{ij},\quad
 I=\sum i h_{ij},\quad J_0=\sum j h_{ij},\\
 \mathcal L_{A,B}=
   \{h:I-Ah_0\in3\mathbb Z_3,\quad
                 J_0-Bh_0\in3\mathbb Z_3\}.
 \end{gathered}
 \tag{E60.2}
\]
Each $\mathcal L_{A,B}$ is a specified index-nine lattice,
free of rank nine. The equations and representatives below
give the homeomorphism in (E60.1).

Let $t$ be tame three-primary inertia and $F$ arithmetic
Frobenius, so $FtF^{-1}=t^4$. A representative section has
values
\[
 \begin{split}
 \nu=v(t)&=c^H,\qquad
 H=\sum h_{ij}e_{ij}+uU+vV,\\
 z=v(F)&=\xi^A\eta^B c^K,
 \qquad
 u=\frac{Ah_0-I}{3},\quad
 v=\frac{Bh_0-J_0}{3}.
 \end{split}
 \tag{E60.3}
\]
The allowed $K$ are exactly the classes in $M/D_{A,B}M$,
where
\[
 \begin{split}
 D_{A,B}e_{ij}&=-3e_{ij}+(4i-A)U+(4j-B)V,\\
 D_{A,B}U&=9U,\qquad D_{A,B}V=9V.
 \end{split}
 \tag{E60.4}
\]
This quotient has the finite coordinates in (E60.1) by
the following actual maps. Put
\[
 p=(\overline{A+1},B),\quad q=(A,\overline{B+1}),
 \quad C_{ij}=\binom{4i-A}{4j-B},
 \quad G=(C_p\ C_q)\in\operatorname{GL}_2(\mathbb Z_3),
 \tag{E60.5}
\]
where the bars denote representatives in $\{0,1,2\}$.
Write $K=\sum k_{ij}e_{ij}+k_UU+k_VV$.
Its seven order-three coordinates are $k_{ij}\bmod3$
for $(i,j)\ne p,q$. Its two order-twenty-seven coordinates are
\[
 \binom{k_p}{k_q}
   +\sum_{(i,j)\ne p,q} k_{ij}G^{-1}C_{ij}
   +3G^{-1}\binom{k_U}{k_V}\pmod {27}.
 \tag{E60.6}
\]
For the inverse, choose the seven residues in $\{0,1,2\}$,
put $k_U=k_V=0$, and choose $k_p,k_q$ by subtracting the
sum in (E60.6) from the two residues in $\{0,\ldots,26\}$.
These are three-adically integral representatives in $M$.

There is also a complete answer to the next central lifting
question within this quotient. Put $R_2=R/\gamma_3R$,
where $\gamma_3R=9\mathbb Z_3Uc\oplus9\mathbb Z_3Vc$.
For a local $R_2$ section, its inertia handle coordinates are
zero. Write its remaining inertia coordinates as in
(E60.3), with $u,v\in\mathbb Z_3/9$ and arbitrary
representatives in $\mathbb Z_3$, and its Frobenius handle
coordinates as $A,B\in\mathbb Z_3$.
The lower tame relation is exactly
\[
 Ah_0-I-3u\in9\mathbb Z_3,\qquad
 Bh_0-J_0-3v\in9\mathbb Z_3.
 \tag{E60.7}
\]
The actual image of local section classes of $R$ consists
exactly of those lower classes satisfying the stronger tests
\[
 \boxed{\quad Ah_0-I-3u\in27\mathbb Z_3,\qquad
               Bh_0-J_0-3v\in27\mathbb Z_3.\quad}
 \tag{E60.8}
\]
The tests are independent of the representatives of $u,v$.
Every nonempty fiber of the map of \emph{full
geometric-conjugacy classes} is a single point.
For a fixed lower representative, the central lift classes
modulo $\gamma_3R$-conjugation form an eighty-one-element
torsor $(\mathbb Z/9)^2$; the lower stabilizer acts simply
transitively on it. This specifies the exact comparison
between the two equivalence relations.

For instance the lower data
$z=1$, $\nu=c^{3U}$ define an $R_2$ section which fails
(E60.8), whereas $z=1$, $\nu=1$ lifts. These two sections
have identical rank-ten abelian data.
All statements concern the actual quotient $R$ of E.58;
the intrinsic source still has the exact kernels in E.58.3.
The manuscript anchor is \S51.16, folios 733v--736v,
scans 172--178, formulas (371)--(379), Part E(c),(d).
\end{statement}

\begin{proof}
The punctured elliptic curve and its three-torsion punctures
have good reduction at two, with the specified integral
tangent. Its pro-three geometric action is unramified there.
The induced action of Frobenius on $R^{\rm ab}/3$ is
unipotent: it is the identity on the two handle and eight
cusp diagonal blocks, with only the handle-to-cusp block
remaining. The cube of this action is the identity.
The kernel of the automorphism action on the Frattini
quotient of a finitely generated pro-three group is pro-three,
as follows from its characteristic exponent-three central
series. Thus the closed Frobenius-action group is pro-three.
Every cocycle consequently factors through the actual
maximal pro-three tame quotient
\[
       \mathcal D=\langle t,F\mid FtF^{-1}=t^4\rangle_3.
 \tag{E60.9}
\]
Here wild inertia is pro-two and contributes no map to the
resulting pro-three semidirect product. The remaining tame
components of prime-to-three order disappear in the same
way. The displayed local quotient and its relation retain
arithmetic Frobenius of residue cardinality four.
Conversely any two values satisfying that relation define
a continuous cocycle, by the presentation of $\mathcal D$.

In the coordinates $(a,b,H)=\xi^a\eta^b c^H$, an exact
group law follows from the commutator module:
\[
 \begin{split}
 &(a,b,H)(A,B,K)=(a+A,b+B,H+K-AbP\\
 &\hspace{19mm}
   +[9Ab(A+1)/2-A\operatorname{aug}H]U\\
 &\hspace{19mm}
   +[9Ab(b+2B+1)/2-B\operatorname{aug}H]V).
 \end{split}
 \tag{E60.10}
\]
For example it follows by moving the commutator module
past $x^{3A}y^{3B}$ and using
$X^3=1+U$, $Y^3=1+V$, $U\mathfrak m=V\mathfrak m=0$.
All its denominators are units in $\mathbb Z_3$.
This law also proves that the thirteen-coordinate
parametrization is a homeomorphism with its product topology.

\emph{Solving the complete tame equation.}
The equation for the two cocycle values is
\[
                       z\varphi(\nu)z^{-1}=\nu^4.
 \tag{E60.11}
\]
On the two handle coordinates it reads
$-2a=4a$, $-2b=4b$. They are torsion-free, so both vanish.
Inertia therefore lies in the abelian module, and we can
write it as in (E60.3). If $z=\xi^A\eta^B c^K$, then
conjugation by $z$ on the module is multiplication by
$X^{3A}Y^{3B}=1+AU+BV$.
Using E58.9 gives the exact error
\[
 z\varphi(\nu)z^{-1}\nu^{-4}
      =c^{4(Ah_0-I-3u)U+4(Bh_0-J_0-3v)V}.
 \tag{E60.12}
\]
This proves (E60.2)--(E60.3), including necessity and
sufficiency. It puts no condition on the Frobenius module
coordinate $K$.

\emph{Every geometric conjugation and the finite coordinates.}
A conjugating element with handle coordinates $(s,t)$
sends the Frobenius handle coordinates to
$(A+3s,B+3t)$, and the inertia central coordinates to
$(u+s h_0,v+t h_0)$. Its nine coordinates $h_{ij}$ are
unchanged. There is therefore a unique choice of handle
coordinates of a conjugator which puts $A,B$ in
$\{0,1,2\}$. After doing so, the remaining conjugators
are precisely $c^L$ with $L\in M$.
They fix the inertia and send the Frobenius coordinate to
\[
      K+(1-AU-BV)L-4\Phi(L)=K+D_{A,B}L.
 \tag{E60.13}
\]
This is (E60.4), with its sign and its full integral domain.

The determinant of $D_{A,B}$ is $(-3)^9\cdot9^2$,
so it has no kernel. To compute its cokernel with actual
representatives, use the two columns in (E60.5).
Modulo three they are the two nonzero coordinate-axis
vectors, so $G$ is invertible over $\mathbb Z_3$.
The corresponding two relations express
\[
       \binom U V=3(G^{\mathsf T})^{-1}\binom{e_p}{e_q}.
 \tag{E60.14}
\]
Every other relation says that
$e_{ij}-(G^{-1}C_{ij})_1e_p-(G^{-1}C_{ij})_2e_q$
has order dividing three. The two remaining relations
$9U=9V=0$ are exactly $27e_p=27e_q=0$ after this
invertible change. No further relations remain.
This proves the complete invariant factors, (E60.6),
and its explicit inverse. In particular the quotient has
order $3^{13}$.

For completeness the rank-nine inertia lattice also has
explicit free coordinates. With $p,q$ as in (E60.5), put
$\alpha_A=\overline{A+1}-A$ and
$\alpha_B=\overline{B+1}-B$, each either $1$ or $-2$.
Choose the seven $h_{ij}$ away from $p,q$ freely, and
choose $r,s\in\mathbb Z_3$. Set
\[
 \begin{split}
 h_p&=\alpha_A^{-1}
       \left(3r-\sum_{(i,j)\ne p,q}(i-A)h_{ij}\right),\\
 h_q&=\alpha_B^{-1}
       \left(3s-\sum_{(i,j)\ne p,q}(j-B)h_{ij}\right).
 \end{split}
 \tag{E60.15}
\]
These solve both congruences and are unique for their
free coordinates. All extraction, division on the closed
principal ideals, residue choices and finite-coordinate
maps are continuous. The explicit construction and its
inverse therefore prove the homeomorphism (E60.1).

\emph{The exact class-three image and stabilizer action.}
In $R_2$, the coordinates $u,v$ are taken modulo nine.
The handle equations still force inertia handles to zero,
and (E60.12) becomes exactly (E60.7).
Changing an inertia lift by $c^{9rU+9sV}$ changes the
error by $c^{-108rU-108sV}$, since Frobenius acts as
$-8$ on $U,V$ and the tame exponent is four.
Thus the error can be killed precisely when its two
parentheses lie in $27\mathbb Z_3$.
Their quotients by $27$ give the unique correction of
the inertia central coordinates. The Frobenius central
coordinates have no further relation, which proves
existence as well as necessity in (E60.8).

The central module $\gamma_3R$ has Frobenius $-8$ and
trivial inertia. An additive cocycle into it has inertia
value killed by $-8-4=-12$, hence zero, and an arbitrary
Frobenius value. Central conjugation changes the latter
by multiplication by $1-(-8)=9$.
The framed lift classes consequently form
$(\gamma_3R)/9(\gamma_3R)\simeq(\mathbb Z/9)^2$.

The lower quotient has a central subgroup generated by
the images of $c^U,c^V$, of order $9^2$.
They are fixed by Frobenius because $-8\equiv1\pmod9$,
and are fixed by every twist of the lower section
because they are central. They are therefore actual
stabilizers of every lower section.
They form its entire lower stabilizer: a fixed Frobenius
value first forces $3s=3t=0$ on the torsion-free handle
coordinates of a conjugator, and then (E60.13) forces
$-3l_{ij}=0$ on each of its nine free module coordinates.
Only the two central coordinates modulo nine remain.
Conjugating a lift by $c^{rU+sV}$ changes its Frobenius
coordinate by
\[
                          9rU+9sV.
 \tag{E60.16}
\]
Modulo $9\gamma_3R=81\mathbb Z_3U+81\mathbb Z_3V$,
this is a bijection from that lower stabilizer to the
framed lift torsor. Hence all framed lifts give the
same full geometric-conjugacy class, proving the
single-point fiber assertion. Finally $\nu=c^{3U}$,
$z=1$ has error $c^{-36U}$, which is zero in $R_2$
and is not removable by a multiple of $108U$.
It gives the stated explicit obstruction.

The exact program checks all nine local Smith systems,
their displayed finite coordinate maps with 108 full
coordinate round trips, 54 tame sections, 108 exact
conjugation and normalization identities, 486 lower
representatives and their lift tests, and all 81
stabilizer values. The proofs above apply to all
three-adic inputs, with these calculations checking
the integral constants and the actual maps.
\end{proof}
